\section{Lean formalization and the exact ProveIt interface}
\label{sec:lean}

\subsection{What was inspected}

The source audit is pinned to ProveIt commit
\texttt{73e6a339db7c72c6b13259059a8345b24385f6e3}.
The inspected project specifies Lean and Mathlib \texttt{v4.32.0}.
The relevant source root is
\begin{center}
\small\texttt{Logic/PrimitiveRecursiveArithmetic/ConciseIPRA/Lean/ConciseIPRA/}.
\end{center}
Its \texttt{PR.lean} defines the actual object theory, and
\texttt{NativeBiInterpretation.lean} provides the translation to the
native scalar calculus~\cite{proveit-pr,proveit-native}.
The packet inventory at
\path{Logic/PrimitiveRecursiveArithmetic/docs/proofs/README.md}
expressly distinguishes delivered research packets from rechecked
proofs~\cite{proveit-packets}. No theorem from those incoming packets is
used as a verified dependency here. The mathematical development above
is self-contained.

Inside the namespace \texttt{LeanProofs.ConciseIPRA}, the function-symbol
syntax is exactly
\begin{lstlisting}
inductive PR : Nat -> Type
  | zero : PR 0
  | succ : PR 1
  | proj {k : Nat} (i : Fin k) : PR k
  | comp {m k : Nat} (f : PR m)
      (gs : Fin m -> PR k) : PR k
  | prec {k : Nat} (f : PR k)
      (g : PR (k + 2)) : PR (k + 1)
\end{lstlisting}
The recursion argument comes first. Terms are \texttt{PTm n}, formulas
are \texttt{Fm PTm n}, the axiom family is \texttt{PRAAx}, and an internal
derivation has type \texttt{Prf PRAAx Gamma phi}. The wrapper for
quantifier-free induction is \texttt{PRA.qfInd}; it requires an explicit
proof that the induction formula satisfies \texttt{Fm.QF}.

\subsection{The three delivered Lean files}

\texttt{NatKernelSpec.lean} gives concrete natural-number definitions of
the coding and checking algorithm. Its functions include
\texttt{cons}, \texttt{head}, \texttt{tail}, \texttt{lookup},
\texttt{zadd}, \texttt{zmul}, dense polynomial addition and convolution,
matrix multiplication, permutation signs, determinants, minors,
adjugates, characteristic coefficients, and finally
\texttt{cayleyResidual} and \texttt{chCheck}. It imports only
\texttt{Std}. For the characteristic polynomial it stores $t$ as one
extra generic variable, then extracts its coefficients. This is
equivalent to the separate $t$-coefficient list used in the manuscript.
Its fixed-cap convolution is multiplication in a truncated polynomial
ring; the support proof in Section~\ref{sec:foundations} is essential
to showing that no relevant coefficient was lost.

These are host-level definitions of numerical functions. The remaining
syntactic step is to construct their \emph{particular} codes in
\texttt{PR k}, with the representation and defining-equation proofs.
The fixed functions passed to folds are implementation abbreviations,
not object-language quantification over all functions.

\texttt{ObjectLanguageTargets.lean} defines the exact target formula.
In schematic ASCII notation, its central definition is
\begin{lstlisting}
def checkStatement (check : PR 1) : Fm PTm 0 :=
  Fm.all (Fm.eq
    (PTm.app check (fun _ : Fin 1 => PTm.var 0))
    (PTm.succ PTm.zero))

def MainDerivationTarget (check : PR 1) : Prop :=
  Prf PRAAx [] (checkStatement check)
\end{lstlisting}
The argument must be the actual code \texttt{chCheckCode : PR 1}
representing the specified numerical checker. The file supplies no
postulated inhabitant of \texttt{MainDerivationTarget chCheckCode}.
Once an inhabitant is constructed, the existing theorem
\texttt{pra\_iff\_native} transfers it to
\texttt{NPrf [] (NFm.ofPRA (checkStatement chCheckCode))}.
That equivalence does not turn arbitrary Lean truth into iPRA provability.

\texttt{HostReference.lean} records ordinary Mathlib comparison targets.
For example,
\begin{lstlisting}
theorem cayleyHamilton (n : Nat)
    (A : Matrix (Fin n) (Fin n) R) :
    aeval A A.charpoly = 0 :=
  Matrix.aeval_self_charpoly A
\end{lstlisting}
Mathlib's theorem and its monic-remainder corollaries are useful for
checking the intended host-level mathematical statement~\cite{mathlib}.
They do not discharge the internal proof task. The three files contain
no added \texttt{axiom} declarations or \texttt{sorry} proof bodies.
Their absence does not mean that the missing target theorem has been
proved: it is deliberately left as a specification.

\subsection{A concrete route to the complete internal proof}

The following order avoids a hidden reflection theorem or an accidental
use of quantified induction.
\begin{enumerate}
\item \textbf{Coding and arithmetic.} Assemble \texttt{PR} syntax for
the finitely many list, signed-integer, comparison, iteration, and
bounded-sum helpers. Derive their equations with
\texttt{PRA.proj}, \texttt{PRA.comp}, \texttt{PRA.precZero}, and
\texttt{PRA.precSucc}. Prove equality tests, finite-sum reindexing, and
two-element cancellation internally.
\item \textbf{Polynomial arrays.} Prove convolution laws coefficient by
coefficient, packed into Boolean equalities. Prove support preservation,
padding compatibility, and specialization of normalized polynomials.
These lemmas give the precise link between the dense checker and the
ordinary integer polynomial notation.
\item \textbf{Matrices and permutations.} Prove the finite matrix ring
laws; implement insertion/deletion of a row and column, transpositions,
and the sign parity calculations. Derive cofactor expansion by the
explicit permutation bijection, then the two adjugate identities.
\item \textbf{Characteristic coefficients.} Construct the finite
coefficient matrices $B_j$. Prove $c_n=1$, $B_{n-1}=I$, and the
coefficient relations used in Theorem~\ref{thm:CH}. These are equalities
of coded arrays, not an application of an evaluation homomorphism whose
commutation hypotheses have not been proved.
\item \textbf{Telescoping.} Apply \texttt{PRA.qfInd} to the packed
version of
$\sum_{j=0}^{k}c_jX^j=-X^{k+1}B_k$ for the bounded prefix.
Handle the final leading coefficient and the empty-matrix branch.
Convert equality of the resulting array to \texttt{chCheck(n)=1},
universally close $n$, and obtain the actual
\texttt{Prf PRAAx [] (checkStatement chCheckCode)}.
\item \textbf{Applications.} Add the explicit monic-division,
Newton/Bell, Gaussian-elimination, interpolation, and finite-tail
invariants already stated in the respective proofs. Keep the input
certificates as assumptions until discharged for a particular
presentation. Apply the native translation only after each internal
derivation has been constructed.
\end{enumerate}

The accompanying \texttt{FORMALIZATION\_PLAN.md} supplies a more detailed
module table, coding equations, finite bounds, and source hashes. This
is substantive remaining formalization work, mainly the internal
finite-array and determinant library. A proof assistant cannot replace
it by observing that the answer of the executable checker is true.

\subsection{Verification status}

No Lean or Lake executable was available in this execution environment.
Consequently neither these outline files nor the upstream project was
compiled here. Their import paths and principal API names were checked
against the pinned source. The delivered PDF was compiled from the
delivered LaTeX, and the executable Python checks were run.

The core verification program computes signed permutation determinants
independently of a characteristic-polynomial library. It checks generic
matrices in dimensions $0$ through $4$ by exact integer-polynomial
normalization, $56$ integer matrices, $180$ matrix instances modulo
$2,4,6,9$, and $30$ rational quaternion representation examples.
It checks both adjugate identities, all coefficient recurrences,
Cayley--Hamilton, polynomial adjugates, and power reduction where
specified. It also verifies the displayed counterexamples and finite
factorial-tail sample bounds.

The separate trace program checks $43$ integer matrices in dimensions
$0$ through $4$ for Newton, Bell, Hessenberg, exact integer
Faddeev--LeVerrier, and truncated exponential/logarithm identities,
as well as the trace obstructions in characteristics $2,3,5$.
These finite checks can detect transcription and implementation errors;
the all-dimension results rest on the human proofs and their induction
accounting, not on extrapolation from the checks.

